\chapter{Audio and Speech Compression}

Marco Cagnazzo - Multimedia Communications University of Padova

\bigskip\hrule\bigskip

\section{Speech and Audio Foundations}

\subsection{Audio in the Multimedia Ecosystem}

\textit{Compression: audio vs. video}

While video accounts for the largest share of network traffic, audio imposes the strictest constraints on the communication chain.

\begin{itemize}
\item \textbf{The Latency Constraint:} human auditory perception is significantly more sensitive to temporal jitter and delay than the visual system. Real-time interaction in immersive environments requires ultra-low end-to-end latency, below 20 ms.
\item \textbf{Dynamic Range and Precision:} audio signals span a dynamic range of up to 100 dB. This necessitates high-precision quantization and adaptive strategies to maintain a constant signal-to-noise ratio.
\item \textbf{Bandwidth vs. Transparency:} the objective is to achieve \textbf{perceptual transparency}, where the compressed signal is indistinguishable from the original, at minimal bitrates for high-density network scenarios.
\end{itemize}

\subsection{Methodological Framework}

\textit{Synergies with Signal Processing Foundations}

The compression of audio signals integrates and specializes the fundamental coding tools developed throughout the course:

\begin{itemize}
\item \textbf{Perceptual Thresholding:} moving from general sound perception to specific \textbf{psychoacoustic masking models} to optimize the distribution of quantization noise.
\item Extending the principles of \textbf{predictive quantization} to the physical modeling of the acoustic source, specifically the human vocal tract.
\item \textbf{Transform Coding:} adapting frequency-domain compaction techniques, such as the DCT, to time-continuous signals via the \textbf{MDCT} (\textit{Modified Discrete Cosine Transform}).
\item \textbf{Entropy Optimization:} applying advanced \textbf{arithmetic coding} or \textbf{Huffman coding} to reach the information-theoretic limit of the quantized representation.
\end{itemize}

\subsection{Methodological Taxonomy}

\textit{Source Modeling vs. Sink Modeling}

Compression strategies are formally categorized based on the specific end of the communication chain they prioritize:

\begin{enumerate}
\item \textbf{Source-Based Coding (Parametric):} exploits the physical constraints of sound generation. Modeling the vocal tract with \textbf{LPC} enables high-intelligibility speech coding at extremely low data rates.
\item \textbf{Sink-Based Coding (Perceptual):} exploits the physiological limitations of the human ear. This is the foundation for high-fidelity audio and music, for example AAC and MP3.
\item \textbf{Hybrid Architectures:} contemporary standards, such as \textbf{Opus} or \textbf{EVS}, dynamically integrate both paradigms to handle mixed content, speech and music, with maximal efficiency.
\end{enumerate}

\begin{calloutbox}{red!70!black}{Emerging Frontier}
\textbf{Statement:} \textbf{Neural codecs} based on generative models are redefining the lower bounds of bitrate by synthesizing high-quality audio from sparse latent representations.

\textbf{Interpretation:} Instead of transmitting a precise waveform description, neural approaches transmit compact latent information that a learned decoder uses to generate perceptually plausible audio.
\end{calloutbox}

\subsection{Temporal Dynamics: Speech vs. Music}

\textit{Variability of the Stationarity Assumption}

Predictive and transform coding rely on signal stationarity. Audio exhibits local stationarity, but the duration of this stability window varies drastically across signal types:

\begin{itemize}
\item \textbf{Speech:} characterized by a quasi-stationarity of approximately 20 ms, dictated by the physical inertia of the human vocal folds and articulators.
\item \textbf{General Audio/Music:} highly heterogeneous.
\begin{itemize}
\item \textit{Tonal segments}, for example a sustained violin note, can remain stationary for hundreds of milliseconds.
\item \textit{Transients}, for example a drum hit, change their statistical properties in less than 2-5 ms.
\end{itemize}
\end{itemize}

\begin{calloutbox}{red!70!black}{Window-Length Compromise}
\textbf{Statement:} modern codecs must switch window lengths dynamically to avoid temporal artifacts like pre-echo during transients while maintaining frequency resolution for tonal parts.

\textbf{Interpretation:} short windows localize attacks in time, while long windows provide better frequency resolution for stable tonal material.
\end{calloutbox}

\subsection{Modeling Paradigms}

\textit{Source Modeling vs. Sink Modeling}

The choice of compression technique depends on the nature of the signal's redundancy:

\begin{itemize}
\item \textbf{Source-Based Modeling (Speech):} effective when the generator is known, namely the vocal tract. We use AR filters to model resonances. It is extremely efficient but fails for non-vocal sounds. The priority is \textit{intelligibility} and low latency.
\item \textbf{Sink-Based Modeling (General Audio):} since music lacks a single physical model, we focus on the receiver, the human auditory system. The priority is \textit{fidelity} and perceptual transparency.
\end{itemize}

\begin{calloutbox}{red!70!black}{Redundancy vs. Irrelevancy}
\textbf{Statement:} speech coding focuses on exploiting \textbf{statistical redundancy} through prediction, while audio coding focuses on removing \textbf{perceptual irrelevancy} through masking.

\textbf{Interpretation:} speech coders model how the signal is produced; perceptual coders model what the listener cannot hear.
\end{calloutbox}

\subsection{Amplitude Statistics across Audio Types}

\textit{PDF Sparsity and Dynamic Range}

The efficiency of optimal quantizers, such as Lloyd-Max quantizers, depends on the shape of the signal's \textbf{probability density function (PDF)}:

\begin{itemize}
\item \textbf{Sparsity:} both speech and many musical instruments exhibit Laplacian-like PDFs, with a sharp peak at zero, indicating that the signal spends most time at low energy.
\item \textbf{Dynamic Range:} music requires 16-24 bit precision to handle the vast range between silence and peaks, whereas speech is often acceptable with 8-bit non-uniform quantization.
\end{itemize}

The following slides provide empirical evidence of these properties using real-world samples.

\subsection{Empirical Analysis: Voiced Speech}

\textit{Vowels: high periodicity and correlation}

\begin{itemize}
\item \textbf{File:} \texttt{vocali.wav}
\item \textbf{Analysis:} the histogram shows a multi-modal distribution reflecting the periodic nature of vowels.
\item \textbf{Correlation:} the plot of consecutive samples, $x_n$ vs. $x_{n-1}$, shows a very tight diagonal alignment, indicating massive redundancy.
\end{itemize}
\begin{center}\includegraphics[width=0.600\linewidth]{Pasted_image_20260513130347.png}\end{center}
\textit{Figura: Histogram of relative frequency of sound samples for vocali.wav — multi-modal distribution around zero}
\begin{center}\includegraphics[width=0.600\linewidth]{Pasted_image_20260513130521.png}\end{center}
\begin{center}\includegraphics[width=0.600\linewidth]{Pasted_image_20260513130558.png}\end{center}

\begin{calloutbox}{gray!60!black}{Figure 2 - Consecutive-sample scatter plot for voiced speech}
\textbf{Placeholder:} add image later from the original PDF. \textbf{Expected content:} log-scale scatter plot of the relative frequency of consecutive samples, showing a tight diagonal alignment that indicates high sample correlation. \textbf{Source cue:} \texttt{[Figura: Scatter plot of relative frequency of consecutive samples (log scale) — tight diagonal alignment indicating high correlation]}
\end{calloutbox}

\subsection{Spectral Analysis: Vowels}

\textit{Formants and the Source-Filter Justification}

\textbf{Visualizing the resonances:}

\begin{itemize}
\item The Fourier analysis of a vowel sound clearly shows the fundamental frequency and harmonics.
\item However, in order to take locality into account, we perform a \textit{short-time Fourier transform} over consecutive time windows.
\item This produces the \textbf{spectrogram}, where we see the \textbf{harmonic structure}: the clear horizontal lines represent the fundamental frequency, or pitch, and its harmonics.
\end{itemize}
\begin{center}\includegraphics[width=0.600\linewidth]{Pasted_image_20260513130738.png}\end{center}

\subsection{Empirical Analysis: General Speech}

\textit{Mixed Phonemes and Laplacian Distribution}

\begin{itemize}
\item \textbf{File:} \texttt{voice.wav}
\item \textbf{Analysis:} a typical speech signal. The histogram exhibits the characteristic Laplacian "peaked" shape centered at zero.
\item \textbf{Coding Insight:} the high density near zero confirms that a large portion of the signal can be represented with very few bits after prediction.
\end{itemize}

\begin{center}\includegraphics[width=0.806\linewidth]{Pasted_image_20260513140029.png}\end{center}

\begin{center}\includegraphics[width=0.806\linewidth]{Pasted_image_20260513140052.png}\end{center}
\subsection{Empirical Analysis: Electronic Music}

\textit{Beats and Transients}

\begin{itemize}
\item \textbf{File:} \texttt{electronic-beat\_10s.wav}
\item \textbf{Analysis:} percussive elements create more dispersion in the consecutive-sample plot.
\item \textbf{Challenge:} sudden transients, the "beats", represent non-stationary events that break the prediction model.
\end{itemize}

\begin{center}\includegraphics[width=0.806\linewidth]{Pasted_image_20260513140114.png}\end{center}

\begin{center}\includegraphics[width=0.806\linewidth]{Pasted_image_20260513140200.png}\end{center}
\subsection{Spectral Analysis: Electronic Music}

\textit{Absence of evident structures}

\begin{itemize}
\item Frequency analysis of music signals does not show evident structures like those observed in speech.
\item The sharp fall of the amplitude spectrum at 16 kHz is due to an acquisition low-pass filter.
\item Using a spectrogram in this case would not improve the analysis.
\end{itemize}
\begin{center}\includegraphics[width=0.806\linewidth]{Pasted_image_20260513140218.png}\end{center}
\subsection{Empirical Analysis: Game Audio}

\textit{Complex Soundscapes and Effects}

\begin{itemize}
\item \textbf{File:} \texttt{battle-of-pogs\_10s.wav}
\item \textbf{Analysis:} the most "noise-like" distribution among the samples.
\item \textbf{Complexity:} high-frequency components and non-harmonic textures make this signal harder to compress using parametric models.
\end{itemize}

\begin{center}\includegraphics[width=0.806\linewidth]{Pasted_image_20260513140307.png}\end{center}
\begin{center}\includegraphics[width=0.806\linewidth]{Pasted_image_20260513140328.png}\end{center}

\subsection{Lessons from Empirical Data}

\textit{Bridging Statistics to Coding Strategy}

\begin{itemize}
\item \textbf{For Speech:} we can rely on source-filter models because the physical constraints of the vocal tract ensure good predictability.
\item \textbf{For Music:} we must shift to perceptual models, because the time-domain signal is too unpredictable for pure source models.
\end{itemize}

\begin{calloutbox}{red!70!black}{Prediction Gain}
\textbf{Statement:} by modeling the resonances seen in the voice examples, we can derive the \textbf{Linear Predictive Coding (LPC)} coefficients to maximize the \textbf{prediction gain} $G_p$.

\textbf{Interpretation:} strong resonance structure makes prediction effective, so the residual carries less variance than the original signal and becomes cheaper to encode.
\end{calloutbox}

\subsection{The Human Speech Production System}

\textit{From Anatomy to Signal Processing}

Speech is generated through a highly structured mechanical process that can be mapped directly to a digital signal processing block diagram:

\begin{itemize}
\item \textbf{The Lungs (Power):} provide the DC air pressure, acting as the system's power supply.
\item \textbf{The Vocal Folds (Source):} the "sound generator." They vibrate periodically for voiced sounds or remain open for noise-like unvoiced sounds.
\item \textbf{The Vocal Tract (Filter):} the pharynx, oral cavity, and nasal cavity. These act as an acoustic tube, shaping the sound through resonances.
\end{itemize}

\subsection{Voiced and Unvoiced Speech}

\textit{The Two Modes of Acoustic Excitation}

The human vocal system generates two fundamentally different types of acoustic excitation based on the state of the glottis:

\begin{itemize}
\item \textbf{Voiced Sounds}, for example vowels like /a/ and /u/, or nasals like /m/:
\begin{itemize}
\item \textbf{Mechanism:} the vocal folds are held together and vibrate periodically due to airflow.
\item \textbf{Signal:} quasi-periodic waveform with a clear \textbf{pitch period} $T_0$.
\item \textbf{Spectrum:} harmonic structure with energy concentrated at multiples of $F_0 = 1/T_0$.
\end{itemize}
\item \textbf{Unvoiced Sounds}, for example fricatives like /s/, /f/, or /sh/:
\begin{itemize}
\item \textbf{Mechanism:} the vocal folds remain open, or abducted; turbulence is created by air rushing through a narrow constriction.
\item \textbf{Signal:} stochastic, noise-like waveform with no apparent periodicity.
\item \textbf{Spectrum:} continuous and "flat" or "colored", depending on the vocal tract shape, similar to white noise.
\end{itemize}
\end{itemize}

\begin{calloutbox}{red!70!black}{Voiced/Unvoiced Coding Decision}
\textbf{Statement:} the encoder must distinguish between these two states: \textbf{voiced} frames require pitch-period estimation, $T_0$, while \textbf{unvoiced} frames are modeled using a random noise generator.

\textbf{Interpretation:} the same vocal-tract filter can be excited by a periodic impulse train for voiced speech or by noise for unvoiced speech.
\end{calloutbox}

\subsection{The Source-Filter Acoustic Model}

\textit{Decoupling Excitation and Resonance}

Fant's classic theory states that the speech signal $s$ is the convolution of an excitation source $e$ and a time-varying vocal tract filter $h$. In the frequency domain:

\[
S(f) = E(f) \cdot H(f)
\]

Instead of encoding the complex waveform $s$, we aim to transmit the filter parameters that represent the vocal tract and a simplified excitation model: an impulse train for voiced speech and white noise for unvoiced speech.

\begin{center}\includegraphics[width=0.600\linewidth]{Pasted_image_20260513140535.png}\end{center}

\subsection{Bridge to Linear Predictive Coding}

\begin{itemize}
\item In the upcoming section, we will learn how to model the $h$ filter by using predictive coding.
\item The model is based on a number $P$ of discrete past samples.
\item If the prediction model is robust, the residual can be easily modeled either as noise or as an impulse train.
\item We work in the discrete domain: the input signal is $x(n)$, and we use the Z-transform to describe the processing architecture.
\item The predictor is a discrete linear filter, which turns out to be related to $H$.
\item \textbf{LPC} means \textbf{Linear Predictive Coding}.
\end{itemize}

\bigskip\hrule\bigskip

\section{Linear Predictive Coding (LPC)}

\subsection{Principles of Linear Prediction}

\textit{Modeling the Signal and the Residual}

Following the source-filter model, we assume the current signal sample $x(n)$ can be estimated from a linear combination of its $P$ past values:

\[
\hat{x}(n) = -\sum_{i=1}^{P} a_i x(n-i)
\]

The prediction error, or \textbf{residual signal} $y(n)$, is defined as:

\[
y(n) = x(n) - \hat{x}(n) = \sum_{i=0}^{P} a_i x(n-i) \quad \text{with } a_0 = 1
\]

\begin{calloutbox}{red!70!black}{Whitening Filter}
\textbf{Statement:} the residual $y(n)$ is the output of a FIR filter $A(z) = \sum_{i=0}^{P} a_i z^{-i}$ applied to $x(n)$. If the prediction is optimal, $y(n)$ is a white-noise-like signal or a periodic impulse train.

\textbf{Interpretation:} LPC attempts to remove the spectral envelope from speech. The residual should contain mainly the excitation information.
\end{calloutbox}

\subsection{Estimation of Optimal Coefficients}

\textit{Minimization of the Residual Variance}

To find the coefficients $\{a_i\}$, we minimize the variance of the residual $\sigma_y^2 = E\{y^2(n)\}$. As seen in the Predictive Quantization module, setting the partial derivatives $\frac{\partial \sigma_y^2}{\partial a_i} = 0$ leads to the \textbf{Yule-Walker equations}:

\[
\mathbf{R}_x \mathbf{a} = -\mathbf{r}_x
\]

where:

\begin{itemize}
\item $\mathbf{R}_x$ is a symmetric \textbf{Toeplitz matrix} of autocorrelations.
\item $\mathbf{r}_x$ is the autocorrelation vector.
\item $\mathbf{a} = [a_1, a_2, \ldots, a_P]^T$ is the vector of LP coefficients.
\end{itemize}

\begin{calloutbox}{red!70!black}{Levinson-Durbin Recursion}
\textbf{Statement:} the Toeplitz structure allows the Yule-Walker system to be solved with $\mathcal{O}(P^2)$ complexity.

\textbf{Interpretation:} this makes real-time estimation feasible even on low-power devices.
\end{calloutbox}

\subsection{The LPC Analysis Process}

\textit{From Waveform to Parametric Representation}

The encoder segments the signal $x(n)$ into frames, approximately 20 ms long, and extracts the main parameters:

\begin{enumerate}
\item \textbf{LP Coefficients} $\{a_i\}$: define the spectral envelope, that is, the vocal tract.
\item \textbf{Gain} $G$: represents the signal energy, $\sigma_y$.
\item \textbf{V/UV Decision:} determines the source type, voiced or unvoiced.
\item \textbf{Pitch Period} $T_0$: fundamental frequency, if voiced.
\end{enumerate}

\begin{center}\includegraphics[width=0.806\linewidth]{Pasted_image_20260513140733.png}\end{center}

\subsection{The Autocorrelation Estimation}

The short-time autocorrelation $R_x(k)$ is the fundamental tool for extracting both source and filter parameters from a windowed frame $x_w(n)$.

\begin{itemize}
\item \textbf{Biased Estimator:} as seen in previous modules, we use:
\end{itemize}

\[
R_x(k) = \sum_{n=0}^{N-1-k} x_w(n) x_w(n+k) \quad \text{for } k = 0, \ldots, K
\]

\begin{itemize}
\item \textbf{Filter Estimation (Lags 0...P):} these values form $\mathbf{R}_x$ and $\mathbf{r}_x$.
\item \textbf{Source Analysis (Long Lags):} the same function is used for:
\begin{itemize}
\item \textbf{Pitch Detection:} finding the maximum peak at $k > P$.
\item \textbf{V/UV Decision:} comparing the peak magnitude against $R_x(0)$.
\end{itemize}
\end{itemize}

\begin{calloutbox}{red!70!black}{Autocorrelation Efficiency}
\textbf{Statement:} by computing $R_x(k)$ once, the encoder determines the spectral envelope, the signal energy, the periodicity, and the excitation type.

\textbf{Interpretation:} one statistic supports both the short-term vocal-tract model and the longer-term pitch analysis.
\end{calloutbox}

\subsection{Characteristics of the Residual $y(n)$}

\textit{What remains after spectral flattening?}

The residual $y(n)$ represents the information that the linear filter cannot capture, namely the excitation source:

\begin{itemize}
\item \textbf{Voiced segments:} $y(n)$ shows sharp, periodic spikes, corresponding to glottal pulses. The prediction gain is high because the signal is highly redundant.
\item \textbf{Unvoiced segments:} $y(n)$ resembles zero-mean white noise. The linear model captures the "color" of the turbulence.
\end{itemize}

\begin{calloutbox}{red!70!black}{LPC Compression Efficiency}
\textbf{Statement:} by transmitting the filter $A(z)$ and a simplified model of $y(n)$, we achieve high compression because $y(n)$ has a much lower variance and flatter spectrum than $x(n)$.

\textbf{Interpretation:} the decoder reconstructs intelligible speech from a compact vocal-tract model plus a compact excitation model.
\end{calloutbox}

\subsection{LPC Synthesis}

\textit{Reconstructing the Signal via Inverse Filtering}

The synthesis process is the exact mathematical inverse of the analysis. If the prediction error filter $A(z)$ is used to extract the residual $y(n)$ from the signal $x(n)$:

\[
Y(z) = A(z)X(z) \implies X(z) = A^{-1}(z)Y(z)
\]

At the decoder, we reconstruct the speech by applying the \textbf{inverse filter} $A^{-1}(z)$ to the excitation:

\begin{itemize}
\item \textbf{Excitation Source:} $y(n)$ is reconstructed as either a periodic impulse train, for voiced speech, or white noise, for unvoiced speech, based on the received gain, pitch, and V/UV parameters.
\item \textbf{Time-Domain Implementation:} the synthesis is performed using a linear recurrence equation:
\end{itemize}

\[
x(n) = y(n) - \sum_{i=1}^{P} a_i x(n-i)
\]

\subsection{LPC Spectral Modeling}

Since $X(z) = A^{-1}(z)Y(z)$, the general shape of $X$ is determined by $A^{-1}(z)$, because $Y$ is either flat, as in white noise, or a Dirac comb, as referenced in page 18 of the source material.

LPC modeling is more accurate near spectral peaks, the \textbf{formants}, than in spectral valleys.

Notice that:

\[
A(z) = a_0 + a_1 z^{-1} + \ldots + a_P z^{-P}
\]

is a FIR filter. Therefore, its inverse is an all-pole IIR filter.

\subsection{The Quantization Problem}

\textit{Why not quantize $\{a_i\}$ directly?}

To transmit the vocal tract model $A(z)$, we need to quantize the filter coefficients $\{a_1, a_2, \ldots, a_P\}$. However, direct scalar quantization of $\{a_i\}$ is extremely dangerous:

\begin{itemize}
\item \textbf{High Dynamic Range:} the coefficients $a_i$ can vary wildly in magnitude, making uniform quantization highly inefficient.
\item \textbf{Lack of Boundedness:} there is no simple relationship to guarantee that the quantized coefficients $\hat{a}_i$ will keep the roots of $A(z)$ inside the unit circle.
\end{itemize}

\begin{calloutbox}{red!70!black}{Stability Constraint}
\textbf{Statement:} if even one pole of the synthesis filter $1/A(z)$ moves outside the unit circle due to quantization noise, the decoder filter becomes unstable, resulting in massive audio distortion.

\textbf{Interpretation:} quantization must preserve filter stability, not only coefficient accuracy.
\end{calloutbox}

\subsection{Line Spectrum Frequencies (LSF)}

\textit{The Itakura Solution}

Introduced by Fumitada Itakura in 1975, \textbf{Line Spectrum Frequencies (LSFs)} map the LPC polynomial $A(z)$ into two new symmetric and anti-symmetric auxiliary polynomials:

\[
P(z) = A(z) + z^{-(P+1)}A(z^{-1})
\]

\[
Q(z) = A(z) - z^{-(P+1)}A(z^{-1})
\]

This decomposition creates a mathematical "miracle" for quantization.

\begin{calloutbox}{red!70!black}{LSF Property 1: Roots on the Unit Circle}
\textbf{Statement:} all roots of $P(z)$ and $Q(z)$ lie exactly on the unit circle. Therefore, they can be represented simply by their angles, or frequencies, $\omega_i \in [0, \pi]$, without needing a magnitude component.

\textbf{Interpretation:} converting LPC coefficients to angles gives a better-behaved parameter space for quantization.
\end{calloutbox}

\subsection{The Interlacing Property}

Let $\omega_i^{(P)}$ be the roots of $P(z)$ and $\omega_i^{(Q)}$ be the roots of $Q(z)$.

\begin{calloutbox}{red!70!black}{LSF Property 2: Strict Interlacing}
\textbf{Statement:} the synthesis filter $1/A(z)$ is strictly stable if and only if the roots of $P(z)$ and $Q(z)$ interlace, or alternate, along the unit circle:

\[
0 < \omega_1^{(P)} < \omega_1^{(Q)} < \omega_2^{(P)} < \omega_2^{(Q)} < \cdots < \pi
\]

\textbf{Interpretation:} stability can be checked and repaired by verifying the order of the LSF angles.
\end{calloutbox}

\textbf{Why this is perfect for coding:}

\begin{itemize}
\item We quantize the frequencies $\omega_i$.
\item If quantization error causes two frequencies to swap order, the decoder detects the anomaly instantly.
\item The decoder simply re-sorts them to restore strict interlacing, guaranteeing a 100\% stable filter.
\end{itemize}

\subsection{From LSFs back to LPC Coefficients}

\textit{Reconstructing the Filter at the Decoder}

Once the decoder receives the quantized LSF angles $\{\omega_i\}$, it reconstructs the filter $A(z)$ to perform synthesis:

\begin{enumerate}
\item \textbf{Form the auxiliary polynomials:} the roots are placed on the unit circle at the received angles:
\end{enumerate}

\[
P(z) = (1 + z^{-1}) \prod_{i=1,3,\ldots}^{P-1} (1 - 2\cos(\omega_i)z^{-1} + z^{-2})
\]

\[
Q(z) = (1 - z^{-1}) \prod_{i=2,4,\ldots}^{P} (1 - 2\cos(\omega_i)z^{-1} + z^{-2})
\]

\begin{enumerate}
\item \textbf{Reconstruct $A(z)$:} based on the initial definitions, we simply average the two polynomials:
\end{enumerate}

\[
A(z) = \frac{P(z) + Q(z)}{2} \tag{1}
\]

\subsection{Visualizing LSFs and Formants}

The LSF representation offers a direct visual link between the root positions and the spectral peaks, or formants, of the vocal tract.

\begin{enumerate}
\item \textbf{Z-Plane (Unit Circle):}
\begin{itemize}
\item Blue roots: $P(z)$.
\item Red roots: $Q(z)$.
\end{itemize}
\item \textbf{Spectral Envelope:} \textbf{resonance} occurs when a $P$-root and a $Q$-root are very close together. In that case, the filter $A(z)$ exhibits a sharp spectral peak, namely a formant.
\end{enumerate}
\begin{center}\includegraphics[width=0.600\linewidth]{Pasted_image_20260513132924.png}\end{center}

\subsection{Efficient Coding: Vector Quantization (VQ)}

Instead of quantizing each LSF $\omega_i$ independently, modern codecs use \textbf{Vector Quantization (VQ)}.

\begin{itemize}
\item \textbf{The Concept:} the entire set of $P$ LSFs is treated as a single vector $\omega = [\omega_1, \ldots, \omega_P]$.
\item \textbf{The Codebook:} encoder and decoder share a "dictionary" of representative LSF vectors, trained on human speech.
\item \textbf{The Search:} the encoder finds the vector in the codebook that minimizes distortion and transmits only its \textbf{index}.
\end{itemize}

\begin{calloutbox}{red!70!black}{LSF Vector Quantization Advantage}
\textbf{Statement:} LSFs are highly correlated because the vocal tract can only take certain physical shapes. VQ exploits this to encode the whole 10th-order filter with just 20-30 bits per frame.

\textbf{Interpretation:} vector quantization exploits dependencies among frequencies better than independent scalar quantization.
\end{calloutbox}

\subsection{LPC-10 Bitrate Breakdown}

The LPC-10 standard, Federal Standard 1015, operates at a constant bitrate of 2.4 kbps, or 54 bits per 22.5 ms frame. The bit allocation adapts to the source type:

\begin{center}
\begin{tabular}{|l|r|r|}
\hline
\textbf{Parameter} & \textbf{Voiced (Bits)} & \textbf{Unvoiced (Bits)} \\ \hline
LPC Coefficients & 41 (Order $P = 10$) & 20 (Order $P = 4$) \\ \hline
Pitch Period ($T_0$) & 7 & 0 (Reserved) \\ \hline
Energy (Gain $G$) & 5 & 5 \\ \hline
Synchronization & 1 & 1 \\ \hline
Error Protection & 0 & 28 \\ \hline
\textbf{Total Bits / Frame} & \textbf{54} & \textbf{54} \\ \hline
\end{tabular}
\end{center}

\begin{calloutbox}{red!70!black}{LPC-10 Bit Allocation Observations}
\textbf{Statement:} voiced frames use maximum spectral detail, $P=10$, and pitch information. Unvoiced frames use lower spectral detail, $P=4$, which leaves significant error-protection bits to improve robustness against channel noise.

\textbf{Interpretation:} the encoder spends bits on the parameters that matter for each excitation type.
\end{calloutbox}

\bigskip\hrule\bigskip
\section{Source-Based Codecs}

\subsection{The CELP Paradigm}

\textit{From Open-Loop to Analysis-by-Synthesis}

LPC-10 failed to achieve high quality because the excitation $y(n)$ was too simplistic. \textbf{CELP} (\textit{Code Excited Linear Prediction}) introduces three innovations:

\begin{enumerate}
\item \textbf{Vector Excitation:} the residual $y(n)$ is no longer just pulses or noise, but a vector chosen from a \textbf{codebook}, or dictionary.
\item \textbf{Analysis-by-Synthesis (AbS):} the encoder acts as a "mini-decoder". It tries different excitation vectors $i$ and different gains $g$, synthesizes the speech, and compares it to the original.
\item \textbf{Closed-Loop Search:} the shape and gain of the vector that produces the minimum error are transmitted.
\end{enumerate}

\begin{center}\includegraphics[width=0.600\linewidth]{Pasted_image_20260513133621.png}\end{center}

\subsection{Perceptual Weighting Theory}

\textit{Hiding Noise in the Formants}

The previous scheme is a simplification made to introduce the main concepts of CELP encoders.

To improve the Quality of Experience, CELP does not minimize the physical error $\varepsilon(n) = x(n) - \hat{x}(n)$. It minimizes a \textbf{perceptually weighted error}.

\begin{itemize}
\item \textbf{The Masking Principle:} the human ear is less sensitive to noise in frequency regions where the signal energy is high, namely the \textit{formants}.
\item \textbf{The Weighting Filter:} CELP uses a filter $W(z)$ to shape the error:
\end{itemize}

\[
W(z) = \frac{A(z)}{A(z/\gamma)} \quad \text{with } 0 < \gamma < 1 \tag{2}
\]

\begin{itemize}
\item \textbf{Pole Shifting:} by multiplying the filter coefficients by $\gamma^i$, we move the poles toward the origin, effectively \textbf{broadening the formants}.
\end{itemize}

\begin{calloutbox}{red!70!black}{CELP Perceptual Weighting}
\textbf{Statement:} CELP minimizes a perceptually weighted error instead of the raw waveform error, using $W(z) = A(z)/A(z/\gamma)$.

\textbf{Interpretation:} the search prefers excitation vectors whose error is hidden under strong speech formants.
\end{calloutbox}

\subsection{Visualizing the Perceptual Mask}

\textit{Spectral Shaping of the Quantization Noise}

\begin{itemize}
\item \textbf{Blue:} the original LPC spectral envelope $1/A(z)$.
\item \textbf{Red:} the broadened envelope $1/A(z/\gamma)$.
\item \textbf{Black:} $W(z)$ de-emphasizes the error near the formants, where the signal provides natural masking, and emphasizes it in the spectral valleys.
\end{itemize}

\begin{center}\includegraphics[width=0.600\linewidth]{Pasted_image_20260513133851.png}\end{center}

\subsection{The CELP Excitation Model (1/2)}

\textit{The Adaptive Codebook and Long-Term Prediction}

The adaptive codebook models the \textbf{periodicity}, or pitch, of voiced speech by "recycling" the previous excitation $y(n)$.

\begin{itemize}
\item \textbf{The Buffer:} a memory of the excitation $y(n)$ synthesized in previous subframes.
\item \textbf{The Lag $Q$ (Pitch Period):} the adaptive vector is extracted by going back $Q$ samples in the buffer:
\end{itemize}

\[
y_{\text{adapt}}(n) = y(n - Q)
\]

\begin{itemize}
\item \textbf{The Gain $b = g_p$ (Pitch Gain):} scales the past excitation to match the current energy.
\end{itemize}

\begin{calloutbox}{red!70!black}{Long-Term Predictor Interpretation}
\textbf{Statement:} the adaptive codebook acts as a \textbf{Long-Term Predictor (LTP)}. It effectively implements a filter $B(z) = g_p z^{-Q}$ that models the harmonic structure of the signal.

\textbf{Interpretation:} pitch periodicity is represented by reusing delayed excitation rather than by storing a new arbitrary vector.
\end{calloutbox}

\subsection{The CELP Excitation Model (2/2)}

\textit{The Fixed Codebook and Multi-Shape Innovation}

Since the past cannot predict everything, such as transitions or unvoiced noise, a \textbf{fixed codebook} provides the "innovation". To achieve high quality without requiring an impossibly large dictionary, the innovation is constructed by combining multiple entries:

\begin{itemize}
\item \textbf{The Structure:} a static dictionary of shape vectors, for example algebraic pulse patterns.
\item \textbf{Multi-Shape Combination:} instead of choosing just one vector, the encoder selects \textbf{two shape vectors}, $v_1$ and $v_2$.
\item \textbf{Independent Gains:} each shape vector is scaled by its own gain factor, $g_{c1}$ and $g_{c2}$.
\end{itemize}

\begin{calloutbox}{red!70!black}{Total CELP Excitation Construction}
\textbf{Statement:} the final excitation $y(n)$ sent to the synthesis filter $1/A(z)$ is the sum of the adaptive past and the multi-shape innovation:

\[
y(n) = \underbrace{g_p \cdot y_{\text{adapt}}(n)}_{\text{Periodic}} + \underbrace{g_{c1} \cdot v_1(n) + g_{c2} \cdot v_2(n)}_{\text{Fixed Innovation}} \tag{3}
\]

\textbf{Interpretation:} the adaptive term represents voiced periodicity, while the fixed-codebook terms represent new unpredictable content.
\end{calloutbox}

\subsection{Global CELP Architecture}

\textit{Implementation in the ITU-T G.729 Standard}
\begin{center}\includegraphics[width=0.600\linewidth]{Pasted_image_20260513134214.png}\end{center}

\begin{itemize}
\item \textbf{Feedback Loop:} the total excitation $y(n)$ is fed back into the adaptive codebook to be used as history for subsequent subframes.
\item \textbf{Subframe Processing:} the excitation parameters, $Q$, $v_1$, $v_2$, $g_p$, $g_{c1}$, and $g_{c2}$, are updated every 5 ms, or subframe, to capture rapid transients, while the LPC filter $A(z)$ is updated every 10 ms, or frame.
\end{itemize}

\subsection{The Analysis-by-Synthesis Loop}

\textit{Filter Memory and the Target Signal ($\hat{p}_0$)}

In the G.729 diagram, before the error is minimized, a signal $\hat{p}_0(n)$ is subtracted from the perceptually weighted input. This is a crucial computational optimization.

\begin{itemize}
\item \textbf{Filter Components:} the output of the synthesis filter $\frac{1}{A(z/\gamma)}$ is the sum of two responses:
\begin{enumerate}
\item \textbf{Zero-Input Response (ZIR):} the "ringing" due to the filter's past memory, namely states from previous subframes.
\item \textbf{Zero-State Response (ZSR):} the response to the \textit{new} excitation, assuming empty memory.
\end{enumerate}
\item \textbf{The Signal $\hat{p}_0(n)$:} this is exactly the ZIR. It is the output the filter would produce if the current codebook excitation were zero.
\end{itemize}

\begin{calloutbox}{red!70!black}{Target Signal in Analysis-by-Synthesis}
\textbf{Statement:} by subtracting the past memory from the original weighted signal, $\text{Target} = x_{\text{weighted}} - \hat{p}_0$, the codebook search is tremendously simplified.

\textbf{Interpretation:} the encoder only needs to match the new codebook vectors, the zero-state response, to this memoryless target signal.
\end{calloutbox}

\subsection{Bitrate Breakdown for G.729}

\textit{8 kbps Allocation: 80 bits per 10 ms frame}

\begin{center}
\begin{tabular}{|l|l|r|}
\hline
\textbf{Parameter} & \textbf{Update Rate} & \textbf{Bits} \\ \hline
LPC Coefficients (LSF) & 10 ms (Once/Frame) & 18 \\ \hline
Adaptive Codebook Index & 5 ms (Twice/Frame) & 14 (8 + 6) \\ \hline
Fixed Codebook Index & 5 ms (Twice/Frame) & 26 (13 + 13) \\ \hline
Codebook Gains & 5 ms (Twice/Frame) & 14 (7 + 7) \\ \hline
Parity/Sync & 10 ms & 8 \\ \hline
\textbf{Total} &  & \textbf{80 bits} \\ \hline
\end{tabular}
\end{center}

\begin{calloutbox}{red!70!black}{G.729 Bitrate Calculation}
\textbf{Statement:} the frame rate and bit allocation give:

\[
R = \frac{80 \text{ bits}}{0.010 \text{ s}} = 8000 \text{ bps}
\]

\textbf{Interpretation:} 80 bits every 10 ms corresponds to an 8 kbps speech codec.
\end{calloutbox}

\subsection{The Wideband Revolution: AMR-WB}

\textit{Standard 3GPP G.722.2: HD Voice}

\textbf{AMR-WB} marked the transition from \textbf{Narrowband (NB)} to \textbf{Wideband (WB)} audio.

\begin{itemize}
\item \textbf{Extended Bandwidth:}
\begin{itemize}
\item \textbf{50-200 Hz:} provides "presence".
\item \textbf{3.4-7 kHz:} improves fricatives and understandability.
\end{itemize}
\item \textbf{ACELP Evolution:} modified perceptual weighting and a significantly larger dictionary.
\item \textbf{Adaptive Rates:} 9 bitrates, from 6.6 to 23.85 kbps.
\end{itemize}

\begin{center}\includegraphics[width=0.600\linewidth]{Pasted_image_20260513134732.png}\end{center}

\subsection{Enhanced Voice Services (EVS)}

\textit{The Standard for VoLTE and 5G}

\textbf{EVS} is the successor to AMR-WB, designed for high-fidelity audio over IP-based packet networks.

\begin{itemize}
\item \textbf{Bandwidth Versatility:} from Narrowband (NB) to \textbf{Full-band (FB, 20 kHz)}.
\item \textbf{Hybrid Coding:} integrates ACELP for speech and MDCT for music modes.
\item \textbf{Bitrate Range:} highly flexible, from 5.9 kbps to 128 kbps.
\end{itemize}

\begin{calloutbox}{red!70!black}{EVS Backwards Compatibility}
\textbf{Statement:} EVS includes an \textbf{AMR-WB IO} (Input/Output) mode to remain compatible with older devices while improving transport robustness.

\textbf{Interpretation:} EVS can serve newer wideband and full-band services while still interoperating with AMR-WB infrastructures.
\end{calloutbox}

\subsection{EVS: Reliability and Quality}

\textit{Solving IP Network Challenges}

EVS is optimized for packet-switched (PS) networks where latency and loss are common:

\begin{itemize}
\item \textbf{Jitter Buffer Management:} new algorithms to smooth out the arrival of IP packets.
\item \textbf{Frame Erasure Concealment (FEC):} advanced techniques to hide packet loss by predicting speech from past parameters.
\item \textbf{Discontinuous Transmission (DTX):} optimized Silence Insertion Descriptors (SID) and comfort noise.
\end{itemize}

\subsection{Mobile Standards Comparison}

\textit{The Path from 2G to 5G}

\begin{center}
\begin{tabular}{|l|l|l|l|}
\hline
\textbf{Standard} & \textbf{Network} & \textbf{Quality} & \textbf{Bandwidth} \\ \hline
GSM-FR & 2G & Telephony & 300-3400 Hz \\ \hline
AMR-NB & 3G & Telephony & 300-3400 Hz \\ \hline
AMR-WB & 4G/LTE & HD Voice & 50-7000 Hz \\ \hline
\textbf{EVS} & \textbf{4G/5G} & \textbf{Hi-Fi} & \textbf{20-20000 Hz} \\ \hline
\end{tabular}
\end{center}

\begin{calloutbox}{red!70!black}{Speech-Audio Convergence}
\textbf{Statement:} EVS represents the final merge of speech-specific ACELP and general-audio MDCT technologies, as developed in the next sections.

\textbf{Interpretation:} modern voice services no longer rely only on source coding; they also need music-quality transform tools.
\end{calloutbox}

\bigskip\hrule\bigskip
\section{Perceptual Audio Coding}

\subsection{Absolute Threshold of Hearing}

\textit{Minimum Audible Pressure}

The human ear does not have a flat frequency response. The power required for a sound to be audible, $S_a(f)$, varies significantly with frequency.

\begin{itemize}
\item \textbf{Audible Range:} typically between 20 Hz and 20 kHz.
\item \textbf{Maximum Sensitivity:} the threshold is lowest, meaning the ear is most sensitive, between 1 kHz and 4 kHz, which corresponds to the core frequencies of human speech.
\item \textbf{Coding Implication:} sounds with power $\sigma^2 < S_a(f)$ are not heard and can be discarded by the encoder without any loss in subjective quality.
\end{itemize}

\begin{center}\includegraphics[width=0.600\linewidth]{Pasted_image_20260513135220.png}\end{center}

\subsection{Critical Bands (BC)}

\textit{The Bark Scale}

The cochlea acts as a bank of overlapping bandpass filters. These filters define the \textbf{critical bands}.

\begin{itemize}
\item \textbf{Definition:} if multiple sinusoids fall within the same critical band, the ear perceives their total power rather than individual components.
\item \textbf{Variable Width:} the bandwidth of these filters increases with frequency: narrower at low frequencies and wider at high frequencies.
\item \textbf{Standardization:} AAC uses 51 groups to approximate these critical bands.
\end{itemize}

\begin{center}\includegraphics[width=0.806\linewidth]{Pasted_image_20260513135313.png}\end{center}

\subsection{Frequency Domain Masking}

\textit{Masking Curves $S_m(f_0, \sigma^2, f)$}

A strong "masker" sound reduces the sensitivity to a second "masked" sound at nearby frequencies.

\begin{itemize}
\item \textbf{Masking Threshold:} for a given frequency $f_0$ and power $\sigma^2$, the threshold $\Phi(f)$ increases, hiding noise below it.
\item \textbf{Asymmetry:} masking is more effective for frequencies higher than the masker's frequency, $f > f_0$.
\item \textbf{Coding Gain:} we allow the quantization noise $S_Q(f)$ to increase as long as it remains below the masking curve: $S_Q(f) < \Phi(f)$.
\end{itemize}

\begin{center}\includegraphics[width=0.600\linewidth]{Pasted_image_20260513135345.png}\end{center}

\subsection{Temporal Domain Masking}

\textit{Pre-masking and Post-masking}

The masking effect also occurs in time, before and after a loud sound.

\begin{itemize}
\item \textbf{Pre-masking:} lasts about 2-5 ms. The ear is "blinded" just before the onset of a loud sound.
\item \textbf{Post-masking:} much longer effect, lasting 100-200 ms. Sensitivity recovers slowly after a sound stops.
\item \textbf{Windowing:} temporal masking justifies the use of specific window shapes to avoid \textbf{pre-echo} artifacts in transients.
\end{itemize}

\begin{center}\includegraphics[width=0.600\linewidth]{Pasted_image_20260513135412.png}\end{center}

\subsection{Perceptual Coding Standards}

\textit{The MPEG Legacy}

The most successful perceptual coders are part of the ISO/IEC MPEG, \textit{Moving Picture Experts Group}, family:

\begin{itemize}
\item \textbf{MP3:} formally known as \textbf{MPEG-1 Audio Layer 3}. It uses a hybrid filterbank, polyphase plus MDCT.
\item \textbf{AAC:} stands for \textbf{Advanced Audio Coding}. It is defined in the \textbf{MPEG-2} and \textbf{MPEG-4} standards.
\end{itemize}

\textbf{Design Objective:} achieve "transparent" quality, indistinguishable from the original, by removing irrelevant information based on the psychoacoustic model.

\subsection{Transform Coding Principles}

\textit{Windowing and Overlap}

Perceptual coders operate on discrete blocks of signal called windows.

\begin{itemize}
\item \textbf{Overlapping:} to avoid blocking artifacts, windows overlap in time.
\item \textbf{Frame vs. Buffer:}
\begin{itemize}
\item \textbf{MP3:} uses a buffer of $N = 512$ samples, coding $M = 32$ new samples.
\item \textbf{AAC:} uses a larger buffer of $N = 2048$ samples, coding $M = 1024$ new samples per frame.
\end{itemize}
\item \textbf{Window Switching:} modern encoders switch between "long" windows, for steady state, and "short" windows, for transients, to manage temporal masking.
\end{itemize}

\subsection{Modified Discrete Cosine Transform (M-DCT)}

\textit{The Engine of Modern Audio Coding}

The \textbf{M-DCT} is a block-based transform used to map the signal from the time domain to the frequency domain.

\begin{itemize}
\item \textbf{Overlap:} it uses overlapping windows, typically 50\%, to eliminate blocking artifacts at frame boundaries.
\item \textbf{TDAC:} \textit{Time Domain Aliasing Cancellation}. The mathematical structure ensures that the aliasing introduced in one block is cancelled by the next during reconstruction.
\item \textbf{Critical Sampling:} this is a key efficiency property.
\begin{itemize}
\item Despite the 50\% overlap, where we process $N = 2M$ samples, the transform produces only $M$ spectral coefficients.
\item \textbf{Result:} the number of frequency coefficients equals the number of \textit{new} time samples, meaning there is \textbf{no redundancy} or data expansion.
\end{itemize}
\end{itemize}

\begin{calloutbox}{red!70!black}{MDCT Critical Sampling}
\textbf{Statement:} with 50\% overlap, MDCT processes $N = 2M$ samples but produces only $M$ spectral coefficients.

\textbf{Interpretation:} overlap improves reconstruction quality and avoids blocking artifacts without increasing the number of transmitted coefficients.
\end{calloutbox}

\subsection{Generic Perceptual Coder Architecture}

\textit{The Analysis-by-Synthesis Loop in Frequency}

The encoder and decoder are asymmetric: most of the complexity lies in the encoder's bit allocation logic.

\textbf{Encoder tasks:}

\begin{enumerate}
\item Time-frequency mapping.
\item Psychoacoustic analysis.
\item Dynamic bit allocation.
\item Quantization and Huffman coding.
\end{enumerate}

\begin{center}\includegraphics[width=0.600\linewidth]{Pasted_image_20260513135624.png}\end{center}

\begin{center}\includegraphics[width=0.600\linewidth]{Pasted_image_20260513135642.png}\end{center}
\subsection{MPEG-1 Layer 3 (MP3): Signal-to-Mask Ratio}

\textit{Signal-to-Mask Ratio (SMR)}

The core efficiency of MP3 comes from its ability to measure the "perceptual headroom" in each subband to guide bit allocation.

\begin{calloutbox}{red!70!black}{Signal-to-Mask Ratio (SMR)}
\textbf{Statement:} for a given frequency band $k$, the SMR is the ratio between the signal power $S_k$ and the masking threshold $\Phi_k$:

\[
SMR_k[\text{dB}] = S_k[\text{dB}] - \Phi_k[\text{dB}]
\]

\textbf{Interpretation:} SMR tells the encoder how much quantization noise can be injected before it becomes audible.
\end{calloutbox}

\begin{itemize}
\item \textbf{Role:} SMR tells the encoder how much noise can be "injected" before it becomes audible.
\item \textbf{Allocation Goal:} assign enough bits to the subband so that the Signal-to-Noise Ratio, $SNR_k$, exceeds $SMR_k$.
\item \textbf{Transparency Condition:} when $SNR_k > SMR_k$, the quantization noise $S_Q$ remains below the threshold $\Phi$, becoming \textbf{inaudible}.
\end{itemize}

\subsection{MPEG-1 Layer 3 (MP3): Psychoacoustic Model}

\textit{Calculating $\Phi(f)$}

To find the global masking threshold $\Phi(f)$, the encoder follows a standardized psychoacoustic procedure:

\begin{enumerate}
\item \textbf{Spectral Analysis:} a high-resolution FFT is performed to identify the energy distribution.
\item \textbf{Identification of Maskers:} the model distinguishes between \textbf{tonal} components, such as sinusoids, and \textbf{noise-like} components, because they have different masking properties.
\item \textbf{Individual Masking Curves:} for each masker, a "spreading function", the masking tent, is calculated based on its intensity and frequency.
\item \textbf{Summation:} all individual curves are combined with the \textbf{Absolute Threshold of Hearing (ATH)} to form the \textbf{Global Masking Threshold} $\Phi(f)$.
\end{enumerate}

\subsection{MPEG-1 MP3 vs. MPEG-2/4 AAC}

\textit{Evolution of Perceptual Audio Coding}

AAC was designed to be the successor to MP3, fixing its structural limitations and providing better quality at lower bitrates.

\begin{center}
\begin{tabular}{|l|l|l|}
\hline
\textbf{Feature} & \textbf{MP3} & \textbf{AAC} \\ \hline
\textbf{Transform} & Hybrid (Polyphase + MDCT) & Pure MDCT \\ \hline
\textbf{Frequency Bands} & 32 uniform subbands & \textbf{51 Critical Bands (SFB)} \\ \hline
\textbf{Resolution} & Fixed / Limited & Flexible, up to 1024 bins \\ \hline
\textbf{Quantization} & Simple / Linear & \textbf{Non-linear, Power Law $X^{3/4}$} \\ \hline
\textbf{Huffman Coding} & Static tables & \textbf{Dynamic, 12 different tables} \\ \hline
\textbf{Typical Bitrate} & Transparent at 192 kbps & \textbf{Transparent at 128 kbps} \\ \hline
\end{tabular}
\end{center}

\begin{calloutbox}{red!70!black}{Scale Factor Logic in AAC}
\textbf{Statement:} unlike MP3's rigid bands, AAC uses the Scale Factor $g$ to surgically place the noise floor for each Scale Factor Band (SFB), ensuring that quantization noise $S_Q$ is always hidden under the masking threshold $\Phi$.

\textbf{Interpretation:} AAC has finer and more flexible control over where quantization noise is placed.
\end{calloutbox}

\subsection{HE-AAC and Unified Codecs}

\textit{Spectral Band Replication (SBR)}

Modern standards like \textbf{HE-AAC}, or High Efficiency AAC, and \textbf{USAC} go beyond simple quantization.

\begin{itemize}
\item \textbf{HE-AAC:} combines AAC with \textbf{SBR} (\textit{Spectral Band Replication}) to replicate high frequencies and with Parametric Stereo.
\item \textbf{USAC:} a hybrid codec that switches between AMR-WB for speech and HE-AAC for music to handle mixed content seamlessly.
\end{itemize}

\bigskip\hrule\bigskip
\section{The Modern Era: Opus and Neural Codecs}

\subsection{Opus: The Hybrid King}

\textit{Blending SILK and CELT}

\textbf{Opus} is designed for the modern web, capable of handling everything from low-bitrate VoIP to high-fidelity music with very low latency.

\begin{itemize}
\item \textbf{SILK Layer:} derived from Skype. Based on \textbf{LPC} (\textit{Linear Prediction}), optimized for speech at 8-12 kbps.
\item \textbf{CELT Layer:} derived from Xiph.org. Based on \textbf{MDCT}, optimized for audio and music at high bitrates.
\item \textbf{The Hybrid Mode:} between 16 and 32 kbps, Opus uses both layers simultaneously, with SILK for the low frequencies and CELT for the highs.
\end{itemize}

\begin{calloutbox}{red!70!black}{Opus Seamless Adaptation}
\textbf{Statement:} Opus can change its bitrate, audio bandwidth, from NB to FB, and coding mode on the fly without any audio glitches or renegotiation.

\textbf{Interpretation:} this makes Opus suitable for unstable real-time internet links where bandwidth and packet loss change continuously.
\end{calloutbox}

\subsection{Opus: Low Latency for the Internet}

\textit{Why it replaced MP3 and AAC for Real-Time}

\begin{itemize}
\item \textbf{Algorithmic Delay:} while MP3 and AAC often require more than 100 ms, Opus can work with frames as small as \textbf{5 ms}.
\item \textbf{VoIP Superiority:} this makes it the standard for WebRTC, Discord, and WhatsApp.
\item \textbf{Resilience:} includes advanced In-band Forward Error Correction (FEC) to survive packet loss on unstable Wi-Fi and 4G connections.
\end{itemize}

\subsection{The AI Frontier: Neural Codecs}

\textit{From Quantization to Generation}

Traditional codecs try to match the \textit{waveform}. Neural codecs, such as Google's \textbf{Lyra} or Meta's \textbf{EnCodec}, try to match the \textit{perception} using generative AI.

\begin{itemize}
\item \textbf{Encoder (Feature Extraction):} a neural network compresses audio into a very small latent space, as low as 3 kbps.
\item \textbf{Decoder (Generative Synthesis):} instead of simple inverse transforms, a \textbf{GAN} (\textit{Generative Adversarial Network}) or a \textbf{diffusion model} re-synthesizes the voice.
\item \textbf{Result:} intelligible and natural-sounding speech at bitrates where CELP would sound like broken noise.
\end{itemize}

\begin{calloutbox}{red!70!black}{Neural Codec Paradigm Shift}
\textbf{Statement:} we no longer send the samples; we send "instructions" for the AI to reconstruct a plausible version of the original signal.

\textbf{Interpretation:} the codec becomes partly generative, so low bitrate is achieved by relying on learned priors rather than only on waveform coding.
\end{calloutbox}

\subsection{The Evolution of Spatial Audio}

\textit{Channels, Objects, and Scenes}

Modern immersive audio, for VR/AR and cinema, uses three different approaches:

\begin{itemize}
\item \textbf{Channel-based:} fixed positions, such as 5.1 or 7.1. The mix is baked for a specific speaker layout.
\item \textbf{Object-based (MPEG-H / Dolby Atmos):} sound sources, for example a "helicopter", are sent as mono signals plus \textbf{XYZ metadata}. The receiver renders the sound based on the user's actual speakers.
\item \textbf{Scene-based (Ambisonics):} captures the entire sound-field sphere. This is ideal for 360 videos and head tracking.
\end{itemize}

\subsection{MPEG-H and Next-Generation TV}

\textit{Interactivity and Personalization}

The MPEG-H standard, used in ATSC 3.0 and DVB, introduces user interactivity:

\begin{itemize}
\item \textbf{Object Control:} the user can manually increase the volume of the "Dialogue Object" or the "Stadium Commentary" while keeping the background noise at the same level.
\item \textbf{Binaural Rendering:} downmixes 3D audio for standard headphones, simulating a 22.2 speaker setup using \textbf{HRTF} (\textit{Head Related Transfer Functions}).
\end{itemize}

\bigskip\hrule\bigskip
\section{Conclusion and Evaluation}

\subsection{Subjective Quality Assessment}

\textit{The Gold Standard: MOS and MUSHRA}

The only true judge of audio quality is the human ear. Standardized tests include:

\begin{itemize}
\item \textbf{MOS (Mean Opinion Score):} listeners rate clips from 1, Bad, to 5, Excellent. A score of 4.0 or higher is usually considered "toll quality".
\item \textbf{MUSHRA (ITU-R BS.1534):} used for intermediate quality. Listeners compare multiple samples against a "Hidden Reference", the original, and "Anchors", low-pass filtered versions.
\item \textbf{Transparency Test:} used for high-end codecs such as AAC. The goal is for the listener to be unable to distinguish the original from the compressed version in a double-blind \textbf{ABX} test.
\end{itemize}

\subsection{Why SNR is Not Effective for Lossy Audio}

\textit{Perceptual Noise vs. Mathematical Noise}

In traditional DSP, we use Signal-to-Noise Ratio (SNR). In audio coding, this fails because:

\begin{itemize}
\item \textbf{Masking:} a codec might introduce 20 dB of noise that is perfectly masked by the signal, therefore inaudible. The SNR would be low, but the quality perfect.
\item \textbf{Phase Shifting:} perceptual codecs often change the phase of the signal. The waveform looks different, with high mathematical error, but it sounds identical.
\end{itemize}

\begin{calloutbox}{red!70!black}{Shift to Perceptual Metrics}
\textbf{Statement:} we need \textbf{perceptual metrics} that model the human auditory system to estimate MOS automatically.

\textbf{Interpretation:} waveform error is not a reliable proxy for human-perceived audio quality.
\end{calloutbox}

\subsection{Objective Metrics: POLQA and ViSQOL}

\textit{Replacing the Human Listener}

Modern networks use algorithms to predict the MOS score in real time.

\begin{itemize}
\item \textbf{POLQA (ITU-T P.863):} the industry standard for 4G/5G. It compares the original and degraded signal using a psychoacoustic model, including Bark scale and loudness mapping.
\item \textbf{ViSQOL (by Google):} a newer metric based on \textbf{spectrogram similarity}. It is particularly robust for VoIP and "generative" neural codecs where the time alignment might be imperfect.
\end{itemize}

\subsection{The Audio Coding Landscape}

\textit{A Comparative Overview}

\begin{center}
\begin{tabular}{|l|l|l|l|l|}
\hline
\textbf{Codec} & \textbf{Paradigm} & \textbf{Bitrate} & \textbf{Latency} & \textbf{Primary Use Case} \\ \hline
AMR-WB & Source (CELP) & 6-24 kbps & 20-30 ms & Legacy Mobile Voice (3G/4G) \\ \hline
AAC & Receiver (MDCT) & 64-128 kbps & \textgreater{}100 ms & Streaming and Storage \\ \hline
\textbf{Opus} & \textbf{Hybrid} & \textbf{6-128 kbps} & \textbf{5-20 ms} & \textbf{WebRTC / Real-time Internet} \\ \hline
EVS & Hybrid & 5-128 kbps & 32 ms & 5G VoNR / VoLTE \\ \hline
Lyra/EnCodec & \textbf{Neural} & \textbf{3-6 kbps} & Variable & Ultra-low bandwidth / Fallback \\ \hline
\end{tabular}
\end{center}

\subsection{Future Trends in Audio Coding}

\textit{The Shift towards AI-Driven and Hybrid Systems}

\begin{itemize}
\item \textbf{The Neural "Safety Net":} today, AI codecs act as a fallback. If the network bitrate drops below 6 kbps, the system switches from waveform coding to generative reconstruction.
\item \textbf{Hybridization:} future standards will likely merge \textbf{MDCT}, for transparency in critical bands, with \textbf{neural layers}, to "imagine" high-frequency textures or background noise.
\item \textbf{Computation vs. Bandwidth:} a paradigm shift where we use more local battery and CPU power, such as NPUs, to save precious network bandwidth.
\item \textbf{The "Hallucination" Challenge:} balancing generative realism with mathematical transparency. In professional and legal fields, faithful reconstruction remains superior to AI-generated "plausible" audio.
\end{itemize}

\section{Final Summary Table}

\begin{center}
\begin{tabular}{|l|l|l|l|}
\hline
\textbf{Concept / Standard} & \textbf{Definition / Formula} & \textbf{Compression role} & \textbf{Notes} \\ \hline
Source-based coding & Parametric coding based on physical sound generation & Achieves very low bitrate for speech & Works best when the source, such as the vocal tract, is known \\ \hline
Sink-based coding & Perceptual coding based on hearing limitations & Removes inaudible or irrelevant information & Main paradigm for music and general audio \\ \hline
Speech stationarity & Speech is quasi-stationary over about 20 ms & Enables frame-based LPC analysis & Music may require longer windows for tones or shorter windows for transients \\ \hline
Source-filter model & $S(f) = E(f) \cdot H(f)$ & Separates excitation from vocal-tract resonances & Voiced excitation is periodic; unvoiced excitation is noise-like \\ \hline
Linear prediction & $\hat{x}(n) = -\sum_{i=1}^{P} a_i x(n-i)$ & Predicts samples from past values and lowers residual variance & Produces a whitening filter $A(z)$ \\ \hline
Residual signal & $y(n) = x(n)-\hat{x}(n)=\sum_{i=0}^{P}a_i x(n-i)$ & Carries excitation information after spectral envelope removal & Periodic for voiced speech, noise-like for unvoiced speech \\ \hline
Yule-Walker equations & $\mathbf{R}_x\mathbf{a}=-\mathbf{r}_x$ & Estimate optimal LPC coefficients & Solved efficiently by Levinson-Durbin recursion \\ \hline
LPC synthesis & $X(z)=A^{-1}(z)Y(z)$ & Reconstructs speech by filtering the excitation & The synthesis filter is all-pole IIR \\ \hline
LSF & Angles obtained from $P(z)$ and $Q(z)$ auxiliary polynomials & Makes LPC coefficient quantization stable and efficient & Roots lie on the unit circle and must interlace \\ \hline
Vector Quantization & Codebook index for an entire LSF vector & Reduces bits for LPC parameters & Exploits correlation among vocal-tract shapes \\ \hline
LPC-10 & 2.4 kbps speech codec with 54 bits per 22.5 ms frame & Very low bitrate speech coding & Low quality because excitation is too simple \\ \hline
CELP & Code Excited Linear Prediction with analysis-by-synthesis & Improves excitation modeling over LPC-10 & Uses adaptive and fixed codebooks \\ \hline
Perceptual weighting & $W(z)=A(z)/A(z/\gamma)$ & Hides CELP error near formants & Broadens formants by pole shifting \\ \hline
G.729 & CELP codec at 8 kbps & Efficient narrowband speech coding & Uses 80 bits per 10 ms frame \\ \hline
AMR-WB & 3GPP G.722.2 wideband speech codec & Extends speech bandwidth to 50-7000 Hz & Known as HD Voice \\ \hline
EVS & Hybrid ACELP/MDCT codec for VoLTE and 5G & Supports NB to full-band audio & Includes packet-loss and jitter robustness tools \\ \hline
Absolute threshold of hearing & Minimum audible power $S_a(f)$ & Allows inaudible components to be discarded & Lowest threshold around 1-4 kHz \\ \hline
Critical bands & Cochlear frequency bands on the Bark scale & Basis for band-wise perceptual bit allocation & AAC approximates them with 51 groups \\ \hline
Frequency masking & A strong masker raises $\Phi(f)$ near its frequency & Allows higher quantization noise where hidden & More effective above the masker's frequency \\ \hline
Temporal masking & Pre-masking lasts 2-5 ms; post-masking lasts 100-200 ms & Guides window switching and pre-echo control & Important for transients \\ \hline
MDCT & Overlapped transform with TDAC and critical sampling & Converts audio into frequency coefficients without expansion & Processes $N=2M$ samples and outputs $M$ coefficients \\ \hline
SMR & $SMR_k[\text{dB}]=S_k[\text{dB}]-\Phi_k[\text{dB}]$ & Determines how many bits a band needs & Transparency requires $SNR_k > SMR_k$ \\ \hline
MP3 & MPEG-1 Audio Layer 3 & Perceptual audio coding with hybrid filterbank & Typically transparent around 192 kbps \\ \hline
AAC & MPEG-2/4 Advanced Audio Coding & More flexible perceptual coding than MP3 & Uses pure MDCT, SFBs, non-linear quantization, and dynamic Huffman tables \\ \hline
HE-AAC / USAC & AAC extensions and unified coding & Improve low-bitrate and mixed-content performance & SBR replicates high frequencies \\ \hline
Opus & Hybrid SILK/CELT codec & Low-latency speech and music over the internet & Supports seamless bitrate, bandwidth, and mode adaptation \\ \hline
Neural codecs & Learned latent coding with generative decoding & Reach very low bitrates by exploiting learned priors & Risk plausible but not mathematically faithful reconstruction \\ \hline
MPEG-H / spatial audio & Channel-, object-, and scene-based immersive audio & Encodes audio scenes and interactive objects & Supports binaural rendering with HRTF \\ \hline
MOS / MUSHRA / ABX & Subjective quality tests & Gold standard for perceived quality & Needed because SNR does not predict audio quality well \\ \hline
POLQA / ViSQOL & Objective perceptual quality metrics & Estimate MOS automatically & Useful for network monitoring and generative codec evaluation \\ \hline
\end{tabular}
\end{center}

